\section{Fields with Potentials and Interactions}
\label{sec:fields_with_potentials_and_interactions}

A gradient term couples neighbouring points, but it does not determine how the field behaves locally when spatial derivatives vanish. Local behaviour is encoded by a potential \(V(\phi)\). A source term describes the influence of a prescribed external distribution, while mixed potentials describe interactions among several fields.

This section derives the resulting nonlinear equations, studies equilibrium and small perturbations, and introduces canonical field momentum and the Hamiltonian density.

\subsection{A general scalar-field density}

Consider

\[
\boxed{
\mathcal{L}
=
\frac{\rho}{2}\phi_{t}^{2}
-
\frac{\kappa}{2}|\nabla\phi|^{2}
-
V(\phi)
+
J(x)\phi
}.
\]

We assume

\[
\rho>0,
\qquad
\kappa>0,
\]

and treat \(J(x)\) as a prescribed source.

The Euler--Lagrange equation is

\[
\boxed{
\rho\phi_{tt}
-
\kappa\nabla^{2}\phi
+
V'(\phi)
=
J
}.
\]

Each term has a clear role:

\begin{itemize}
    \item \(\rho\phi_{tt}\) is the inertial response;
    \item \(-\kappa\nabla^{2}\phi\) is the spatial restoring contribution;
    \item \(V'(\phi)\) is the local force generated by the potential;
    \item \(J\) drives the field externally.
\end{itemize}

The equation is nonlinear whenever \(V'(\phi)\) is nonlinear in \(\phi\).

\subsection{Derivation term by term}

The derivative with respect to the time derivative is

\[
\frac{\partial\mathcal{L}}{\partial\phi_{t}}
=
\rho\phi_{t}.
\]

Therefore

\[
\frac{\partial}{\partial t}
\left(
\frac{\partial\mathcal{L}}{\partial\phi_{t}}
\right)
=
\rho\phi_{tt}.
\]

For the spatial derivatives,

\[
\frac{\partial\mathcal{L}}
{\partial(\partial_{i}\phi)}
=
-\kappa\partial_{i}\phi,
\]

so

\[
\sum_{i}
\partial_{i}
\left(
\frac{\partial\mathcal{L}}
{\partial(\partial_{i}\phi)}
\right)
=
-\kappa\nabla^{2}\phi.
\]

Finally,

\[
\frac{\partial\mathcal{L}}{\partial\phi}
=
-V'(\phi)+J.
\]

Substitution into the field Euler--Lagrange equation gives the stated result.

\subsection{Uniform equilibrium fields}

A uniform static field satisfies

\[
\phi_{t}=0,
\qquad
\nabla\phi=0.
\]

The field equation reduces to

\[
\boxed{
V'(\phi_{0})=J
}.
\]

In the source-free case,

\[
J=0,
\]

uniform equilibria are stationary points of the potential:

\[
V'(\phi_{0})=0.
\]

The second derivative distinguishes their local character:

\[
V''(\phi_{0})>0
\]

indicates a local minimum and usually stable small oscillations, while

\[
V''(\phi_{0})<0
\]

indicates a local maximum and instability.

\subsection{Quadratic potential}

Let

\[
V(\phi)
=
\frac{m^{2}}{2}\phi^{2}.
\]

Then

\[
V'(\phi)=m^{2}\phi.
\]

The field equation is

\[
\boxed{
\rho\phi_{tt}
-
\kappa\nabla^{2}\phi
+
m^{2}\phi
=
J
}.
\]

When \(J=0\), this is a linear wave equation with a local restoring term. The parameter \(m^{2}\) controls the strength of the restoring force.

For a spatially uniform field,

\[
\nabla\phi=0,
\]

the equation becomes

\[
\rho\phi_{tt}+m^{2}\phi=0.
\]

Thus every spatial point behaves locally like a harmonic oscillator when gradient coupling is absent.

\subsection{Plane-wave check for the quadratic theory}

For \(J=0\), try

\[
\phi(\mathbf{x},t)
=
A e^{i(\mathbf{k}\cdot\mathbf{x}-\omega t)}.
\]

Then

\[
\phi_{tt}=-\omega^{2}\phi,
\qquad
\nabla^{2}\phi=-|\mathbf{k}|^{2}\phi.
\]

Substitution gives

\[
-\rho\omega^{2}
+
\kappa|\mathbf{k}|^{2}
+
m^{2}
=0.
\]

Therefore

\[
\boxed{
\omega^{2}
=
\frac{\kappa}{\rho}|\mathbf{k}|^{2}
+
\frac{m^{2}}{\rho}
}.
\]

At zero wave number,

\[
\omega^{2}(0)=\frac{m^{2}}{\rho}.
\]

The field can therefore oscillate even when it is spatially uniform.

\subsection{A quartic interaction}

Consider

\[
V(\phi)
=
\frac{m^{2}}{2}\phi^{2}
+
\frac{\lambda}{4}\phi^{4},
\]

with \(\lambda>0\). Then

\[
V'(\phi)
=
m^{2}\phi
+
\lambda\phi^{3}.
\]

The source-free field equation is

\[
\boxed{
\rho\phi_{tt}
-
\kappa\nabla^{2}\phi
+
m^{2}\phi
+
\lambda\phi^{3}
=0
}.
\]

The cubic term in the equation makes the theory nonlinear. If \(\phi_{1}\) and \(\phi_{2}\) are solutions, their sum is generally not a solution.

The quartic term also makes the potential grow positively for large \(|\phi|\), so the local potential energy is bounded below when \(\lambda>0\).

\subsection{Double-well potential}

A useful nonlinear example is

\[
\boxed{
V(\phi)
=
\frac{\lambda}{4}
(\phi^{2}-v^{2})^{2}
},
\]

where \(\lambda>0\) and \(v>0\).

Its derivative is

\[
V'(\phi)
=
\lambda\phi(\phi^{2}-v^{2}).
\]

The source-free field equation is

\[
\boxed{
\rho\phi_{tt}
-
\kappa\nabla^{2}\phi
+
\lambda\phi(\phi^{2}-v^{2})
=0
}.
\]

Uniform equilibria satisfy

\[
\lambda\phi_{0}(\phi_{0}^{2}-v^{2})=0.
\]

Therefore

\[
\phi_{0}=0,
\qquad
\phi_{0}=v,
\qquad
\phi_{0}=-v.
\]

The second derivative is

\[
V''(\phi)
=
\lambda(3\phi^{2}-v^{2}).
\]

At \(\phi_{0}=0\),

\[
V''(0)=-\lambda v^{2}<0,
\]

so the central equilibrium is unstable. At \(\phi_{0}=\pm v\),

\[
V''(\pm v)=2\lambda v^{2}>0,
\]

so the two outer equilibria are stable against sufficiently small uniform perturbations.

\subsection{Linearization around an equilibrium}

Let \(\phi_{0}\) be a constant equilibrium satisfying

\[
V'(\phi_{0})=J_{0}.
\]

Write

\[
\phi(x)=\phi_{0}+\psi(x)
\]

and

\[
J(x)=J_{0}+\delta J(x),
\]

where \(\psi\) and \(\delta J\) are small.

Expand the potential derivative:

\[
V'(\phi_{0}+\psi)
=
V'(\phi_{0})
+
V''(\phi_{0})\psi
+
O(\psi^{2}).
\]

Substitution into the field equation gives

\[
\rho\psi_{tt}
-
\kappa\nabla^{2}\psi
+
V''(\phi_{0})\psi
=
\delta J
+
O(\psi^{2}).
\]

Keeping only first-order terms,

\[
\boxed{
\rho\psi_{tt}
-
\kappa\nabla^{2}\psi
+
V''(\phi_{0})\psi
=
\delta J
}.
\]

Thus small perturbations around a nonlinear equilibrium obey an approximately linear field equation. The curvature \(V''(\phi_{0})\) determines the local restoring strength.

\subsection{Stability from the linearized equation}

For a source-free spatially uniform perturbation,

\[
\rho\psi_{tt}
+
V''(\phi_{0})\psi
=0.
\]

If

\[
V''(\phi_{0})>0,
\]

the perturbation oscillates with frequency

\[
\omega_{0}^{2}
=
\frac{V''(\phi_{0})}{\rho}.
\]

If

\[
V''(\phi_{0})<0,
\]

the equation has exponentially growing solutions. The equilibrium is linearly unstable.

This connects the geometry of the potential with dynamical stability.

\subsection{Interactions among several fields}

For two fields, a general local potential is

\[
V=V(\phi,\chi).
\]

The density

\[
\mathcal{L}
=
\frac{1}{2}\phi_{t}^{2}
-
\frac{c_{\phi}^{2}}{2}|\nabla\phi|^{2}
+
\frac{1}{2}\chi_{t}^{2}
-
\frac{c_{\chi}^{2}}{2}|\nabla\chi|^{2}
-
V(\phi,\chi)
\]

produces

\[
\boxed{
\phi_{tt}
-
c_{\phi}^{2}\nabla^{2}\phi
+
\frac{\partial V}{\partial\phi}
=0
},
\]

and

\[
\boxed{
\chi_{tt}
-
c_{\chi}^{2}\nabla^{2}\chi
+
\frac{\partial V}{\partial\chi}
=0
}.
\]

A mixed potential such as

\[
V_{\mathrm{int}}
=
g\phi^{2}\chi^{2}
\]

contributes

\[
\frac{\partial V_{\mathrm{int}}}{\partial\phi}
=
2g\phi\chi^{2},
\]

and

\[
\frac{\partial V_{\mathrm{int}}}{\partial\chi}
=
2g\phi^{2}\chi.
\]

Each field changes the effective local force acting on the other.

\subsection{Energy density}

For the density

\[
\mathcal{L}
=
\frac{\rho}{2}\phi_{t}^{2}
-
\frac{\kappa}{2}|\nabla\phi|^{2}
-
V(\phi)
+
J\phi,
\]

the corresponding energy density is

\[
\boxed{
\mathcal{E}
=
\frac{\rho}{2}\phi_{t}^{2}
+
\frac{\kappa}{2}|\nabla\phi|^{2}
+
V(\phi)
-
J\phi
}.
\]

The kinetic and gradient contributions are positive when \(\rho>0\) and \(\kappa>0\). Stability also depends on whether the effective potential

\[
V_{\mathrm{eff}}(\phi,x)
=
V(\phi)-J(x)\phi
\]

is bounded below.

If the source depends explicitly on time, the field can exchange energy with the external system, and the field energy need not be conserved.

\subsection{Canonical field momentum}

In mechanics, the momentum conjugate to \(q\) is

\[
p=\frac{\partial L}{\partial\dot q}.
\]

For a field, the canonical momentum density conjugate to \(\phi\) is defined by

\[
\boxed{
\pi(\mathbf{x},t)
:=
\frac{\partial\mathcal{L}}
{\partial\phi_{t}}
}.
\]

For the general scalar density,

\[
\pi
=
\rho\phi_{t}.
\]

Thus the canonical momentum is itself a field. It assigns a momentum density to every spatial point.

For several fields,

\[
\pi_{a}
=
\frac{\partial\mathcal{L}}
{\partial(\partial_{t}\phi^{a})}.
\]

\subsection{Hamiltonian density}

The Hamiltonian density is defined by the field-theoretic Legendre transform

\[
\boxed{
\mathcal{H}
=
\pi\phi_{t}-\mathcal{L}
}.
\]

For

\[
\pi=\rho\phi_{t},
\]

we obtain

\[
\begin{aligned}
\mathcal{H}
&=
\rho\phi_{t}^{2}
-
\left[
\frac{\rho}{2}\phi_{t}^{2}
-
\frac{\kappa}{2}|\nabla\phi|^{2}
-
V(\phi)
+
J\phi
\right]\\
&=
\frac{\rho}{2}\phi_{t}^{2}
+
\frac{\kappa}{2}|\nabla\phi|^{2}
+
V(\phi)
-
J\phi.
\end{aligned}
\]

Therefore

\[
\boxed{
\mathcal{H}
=
\frac{\pi^{2}}{2\rho}
+
\frac{\kappa}{2}|\nabla\phi|^{2}
+
V(\phi)
-
J\phi
}.
\]

For this standard theory, the Hamiltonian density equals the energy density.

The Hamiltonian functional is

\[
\boxed{
H[\phi,\pi]
=
\int_{\Sigma}
\mathcal{H}
\,\mathrm{d}^{3}x
}.
\]

A full Hamiltonian formulation will not be developed here. The important introductory point is that both the field value and its conjugate momentum are spatially distributed variables.

\subsection{When the Legendre transform fails}

The definition of \(\mathcal{H}\) requires that the relation

\[
\pi
=
\frac{\partial\mathcal{L}}{\partial\phi_{t}}
\]

can be inverted to express \(\phi_{t}\) in terms of \(\phi\) and \(\pi\).

For

\[
\pi=\rho\phi_{t}
\]

with \(\rho\neq0\), inversion is immediate:

\[
\phi_{t}=\frac{\pi}{\rho}.
\]

In theories with constraints or gauge redundancy, this inversion may fail. The Hamiltonian formulation then requires a separate constraint analysis. Electromagnetism will later provide an important example.

\subsection{Diagnostic checks}

For a theory with a potential or interaction, one should ask:

\begin{enumerate}
    \item Is the potential differentiated with the correct sign?
    \item Is the source dynamical or prescribed?
    \item Is the potential bounded below?
    \item Which constant fields satisfy the equilibrium equation?
    \item What is the sign of \(V''\) at each equilibrium?
    \item Does linearization reproduce the expected small-oscillation equation?
    \item Is the canonical momentum definition invertible?
    \item Does the Hamiltonian density have the expected energy signs?
\end{enumerate}

\subsection{Common mistakes}

Typical errors include:

\begin{enumerate}
    \item writing \(V(\phi)\) rather than \(V'(\phi)\) in the field equation;
    \item changing the sign of the potential contribution incorrectly;
    \item applying superposition to a nonlinear equation;
    \item calling every stationary point stable;
    \item linearizing without stating which terms are neglected;
    \item varying an external source that should remain fixed;
    \item identifying \(\mathcal{L}\) with the energy density;
    \item forgetting that canonical momentum is a density field.
\end{enumerate}

\subsection{Chapter summary}

This chapter extended the action principle from particle trajectories to fields. The field action is

\[
S[\phi]
=
\int_{\Omega}
\mathcal{L}(\phi,\partial_{\mu}\phi,x)
\,\mathrm{d}^{4}x.
\]

A varied field is written as

\[
\phi_{\varepsilon}
=
\phi+
\varepsilon\eta,
\]

and its derivative varies according to

\[
\delta(\partial_{\mu}\phi)
=
\partial_{\mu}(\delta\phi).
\]

Spacetime integration by parts separates the first variation into an interior term and a boundary term. For fixed boundary values, stationarity gives

\[
\boxed{
\partial_{\mu}
\left(
\frac{\partial\mathcal{L}}
{\partial(\partial_{\mu}\phi)}
\right)
-
\frac{\partial\mathcal{L}}{\partial\phi}
=0
}.
\]

The same method produced coupled equations for several fields, the wave equation for a vibrating string, and the Laplace and Poisson equations from a static energy functional. Local potentials produced nonlinear equations, equilibrium conditions, and linearized small-perturbation theories.

Finally, the canonical momentum density and Hamiltonian density were identified as

\[
\pi
=
\frac{\partial\mathcal{L}}
{\partial\phi_{t}},
\qquad
\mathcal{H}
=
\pi\phi_{t}-\mathcal{L}.
\]

The practical objective of the chapter has therefore been achieved: given a Lagrangian density, one can define admissible variations, control the boundary term, derive the field equation, and interpret the roles of kinetic, gradient, source, potential, and interaction terms.

The next chapter introduces Minkowski spacetime and Lorentz-covariant notation, allowing these field equations to be written in a form adapted to relativistic symmetry.